\subsection{关于点测度积分的、取值于 Banach 空间函数的积分}

定理 2. --- 设 \((\pi, g)\) 是一个 \(\mu\) -适应对，并令

\[
\nu = \int g (t) \varepsilon_ {\pi (t)} d \mu (t).
\]

设 \(\mathbf{f}\) 是定义在 \(\mathbf{X}\) 上的一个函数，取值于 Banach 空间 \(\mathbf{F}\) 或 \(\overline{\mathbf{R}}\) 中。\(\mathbf{f}\) 为本质 \(\nu\) -可积的，必须且只须 \(t \mapsto \mathbf{f}(\pi(t))g(t)\) 是本质 \(\mu\) -可积的，此时

\[
\int \mathbf {f} (x) d \nu (x) = \int \mathbf {f} (\pi (t)) g (t) d \mu (t).\tag{9}
\]

此外假设 \(\pi\) 是连续且逆紧的，g 是连续的且使 \(g(t) > 0\) 对所有 \(t \in T\) 成立。那么，f 为 \(\nu\) -可积的，必须且只须 \(t \mapsto \mathbf{f}(\pi(t))g(t)\) 是 \(\mu\) -可积的。

\begin{enumerate}
\def\labelenumi{\Alph{enumi})}
\tightlist
\item
  我们首先处理测度 \(\mu\) 具有紧支集 K 且 g 在其上有界的情形。此时测度 \(\mu\) 和 \(\nu\) 都是有界的，因而可以把陈述中的“本质可积”换成“可积”。假设 f 是 \(\nu\) -可积的：则函数 \(\mathbf{f}(\pi(t))g(t)\) 是 \(\mu\) -可积的，且依据 §3, No.~3 的定理 1，关系式 (9) 成立。反之，假设 \(\mathbf{f}(\pi(t))g(t)\) 是 \(\mu\) -可积的：则 f 是 \(\nu\) -可测的（No.~3, 命题 3 及注），并且
\end{enumerate}

\[
\int^ {\bullet} | \mathbf {f} (x) | d \nu (x) = \int^ {\bullet} | \mathbf {f} (\pi (t)) | g (t) d \mu (t) <   + \infty
\]

（No.~2, 定理 1）；因此 \(\mathbf{f}\) 是本质 \(\nu\) -可积的（§1, No.~3, 命题 9），从而是 \(\nu\) -可积的。于是 §3, No.~3 的定理 1 蕴涵 (9)。

\begin{enumerate}
\def\labelenumi{\Alph{enumi})}
\setcounter{enumi}{1}
\tightlist
\item
  现在转到一般情形。设 \(\mathfrak{K}\) 是由紧子集 \(K\) （\(T\) 中的、使 \(g|K\) 连续者）组成的集：\(\mathfrak{K}\) 是 \(\mu\) -稠密的（Ch. IV, §5, No.~10, 命题 15），因此测度 \(\mu\) 是一个测度族 \((\mu_{\alpha})_{\alpha \in A}\) 之和，这些测度的支集都是 \(\mathfrak{K}\) 的元素（§2, No.~3, 命题 4）。对 \((g,\pi)\) 显然是 \(\mu_{\alpha}\) -适应的（对每个 \(\alpha \in A\) ），且测度 \(\nu\) 是测度族 \(\nu_{\alpha} = \int \varepsilon_{\pi(t)} g(t) d\mu_{\alpha}(t)\) 之和（§3, No.~1, 命题 1 的推论）。由于 A) 的论证可以应用于测度 \(\mu_{\alpha}, \nu_{\alpha}\) ，陈述的第一部分便由 §2, No.~2 的命题 3 得出。
\end{enumerate}

函数 \(\mathbf{f}\) （相应地，\(t\mapsto \mathbf{f}(\pi (t))g(t)\) ）为关于 \(\nu\) （相应地，关于 \(\mu\) ）可积的，必须且只须它是本质可积的，并且

\[
\int^ {*} | \mathbf {f} (x) | d \nu (x) <   + \infty \quad (\mathrm{resp.} \int^ {*} | \mathbf {f} (\pi (t)) | g (t) d \mu (t) <   + \infty).
\]

因此陈述的第二部分由第一部分及命题 2 得出。

注. --- 设 \((\pi, g)\) 是一个 \(\mu\) -适应对，\(\pi'\) 是 T 到 X 中的一个映射，\(g'\) 是定义在 T 上的一个有限数值函数 \(\geqslant 0\) ，使得 \(\pi'\) （相应地，\(g'\) ）等于 \(\pi\) （相应地，g）局部几乎处处（关于 \(\mu\) ）。则对 \((\pi', g')\) 是 \(\mu\) -适应的，测度 \(\lambda_t = g(t)\varepsilon_{\pi(t)}\) 与 \(\lambda_t' = g'(t)\varepsilon_{\pi'(t)}\) 局部几乎处处相等，且 \(\int g(t)\varepsilon_{\pi(t)} d\mu(t) = \int g'(t)\varepsilon_{\pi'(t)} d\mu(t)\) 。现在，如果 \(\pi'\) 和 \(g'\) 只是局部几乎处处（关于 \(\mu\) ）有定义，并且存在一个 \(\mu\) -适应对 \((\pi, g)\) 使得 \(\pi'\) （相应地，\(g'\) ）局部几乎处处等于 \(\pi\) （相应地，g），我们仍称对 \((\pi', g')\) 是 \(\mu\) -适应的，并置

\[
\int g ^ {\prime} (t) \varepsilon_ {\pi^ {\prime} (t)} d \mu (t) = \int g (t) \varepsilon_ {\pi (t)} d \mu (t)
\]

（参见 §3, No.~3 的注）。当 \(\pi\) 和 \(g\) 仅被假设为局部几乎处处有定义时，定理 1、定理 2 及命题 3 的陈述仍然有效。

